\section{Join-Size Bounds as Moment Problems}
\label{sec:theory}

\begin{figure}[t]
\centering
\begin{tikzpicture}[
  font=\footnotesize,
  node distance = 4mm and 5mm,
  box/.style = {draw, rounded corners=1pt, align=center,
                inner sep=3pt, fill=blue!4},
  lab/.style = {font=\scriptsize\itshape, text=black!60},
  arr/.style = {-{Stealth[length=2mm]}, thick, black!70}]

\node[box] (field) {degree field $\mathcal{D}$\\[1pt]
  $d(x_1){=}(8,9)$\; $d(x_2){=}(4,3)$\; $d(x_3){=}(2,1)$};
\node[lab, above=1mm of field] {per-key degree vectors};

\node[box, right=of field, yshift=7mm] (mom) {observed moments\\
  $\mu_\alpha = \sum_x d(x)^\alpha$};
\node[box, right=of field, yshift=-7mm] (omega) {worst-case domain\\
  $\Omega = \prod_i \mathrm{supp}_i$};

\node[box, right=of mom, yshift=-7mm, text width=24mm] (clp)
  {certificate LP\\ $\min_c \sum_\alpha c_\alpha \mu_\alpha$\\
   s.t.\ $\sum_\alpha c_\alpha d^\alpha \ge \prod_i d_i$};

\node[box, right=of clp] (bound) {bound $B$\\ $+$ proof\\certificate};

\node[box, below=7mm of clp, text width=30mm] (elp)
  {entropic LP (equivalent)\\ support constraints
   $h(\cup_{i\in S}A_i) \le \log\mu_\alpha$};

\draw[arr] (field) -- (mom);
\draw[arr] (field) -- (omega);
\draw[arr] (mom) -- (clp);
\draw[arr] (omega) -- (clp);
\draw[arr] (clp) -- (bound);
\draw[arr, dashed] (mom) |- (elp);
\draw[arr, dashed] (elp) -- (bound);
\end{tikzpicture}
\caption{The \framework{} pipeline.  A join induces a degree field on
its separator; statistics are moments of the field; a bound is either
a polynomial certificate over the worst-case domain (top path) or a
support constraint inside the entropic LP (bottom path) --- two views
of the same moment problem.}
\label{fig:pipeline}
\end{figure}

We now re-derive join-size bounds from scratch --- no entropies, no
$h(U)$ variables --- as a moment problem over a combinatorial object
we call the \emph{degree field}.  All prior bounds reappear as
sections of one linear program.

\subsection{The Degree Field and the Product Moment}

\begin{definition}[Degree field]
\label{def:field}
Let $Q$ be a \emph{star query} on separator $X$, i.e., $A_i \cap A_j =
\{X\}$ for $i\neq j$, or more generally a query decomposed at a
separator attribute $X$.  For each key $x \in \mathrm{dom}(X)$ its
\emph{degree vector} is
$d(x) = \big(d_1(x),\dots,d_m(x)\big) \in \mathbf{N}^m$ with
$d_i(x) = \big|\sigma_{X=x}(R_i)\big|$
(the tuple multiplicity, including row multiplicity under bag
semantics).  The multiset $\mathcal{D} = \{d(x)\}_x$ is the
\emph{degree field} of $Q$ at $X$.
\end{definition}

\begin{proposition}[Product moment]
\label{prop:prod}
For a star query on $X$, $|Q| = \sum_x \prod_i d_i(x)$.  For a general
acyclic query with join tree $\mathcal{T}$, $|Q|$ is a product of such
sums over the separators of $\mathcal{T}$.
\end{proposition}

\begin{definition}[Observed moments]
A \emph{statistic} is a moment
$\mu_\alpha = \sum_x d(x)^\alpha$, $\alpha \in \mathbf{N}^m$.
\end{definition}

\noindent Familiar statistics are moments of particular orders:

\begin{center}\small
\begin{tabular}{c|c|l}
$\alpha$ & moment $\mu_\alpha$ & known statistic \\ \hline
$p\cdot e_i$ & $\sum_x d_i(x)^p$ & $\|\deg_i\|_p^p$ \;(all of LpBound) \\
$e_i{+}e_j$ & $\sum_x d_i d_j$ & $|R_i \Join_X R_j|$ \;(pair count) \\
$e_i{+}e_j{+}e_k$ & $\sum_x d_i d_j d_k$ & triple count \\
$\mathbf{0}$ & $\sum_x 1$ & $|\mathrm{dom}(X)|$ \;(support) \\
\end{tabular}
\end{center}

\noindent Note what this view explains immediately: univariate moments
$\{p\,e_i\}$ are \emph{permutation-blind} within each axis --- they
cannot see which $d_i(x)$ pairs with which $d_j(x)$.  The join size is
a \emph{joint} moment; observing it (or any $\alpha$ with $\ge 2$
nonzero coordinates) is precisely the fix.

\subsection{The Certificate LP}

Fix a moment set $\mathcal{A} \subset \mathbf{N}^m$ and let
$\boldsymbol{\mu} = (\mu_\alpha)_{\alpha\in\mathcal{A}}$ be observed
on the realized field.  Since scalar moments do not reveal the
alignment of coordinates, a valid certificate must dominate the
product on the \emph{worst-case domain}
\begin{equation}
\Omega \;=\; \prod_{i=1}^{m} \big( \{0\} \cup \mathrm{supp}(\deg_i)\big)
\;\subseteq\; \mathbf{N}^m .
\label{eq:omega}
\end{equation}

\begin{definition}[Certificate bound]
\label{def:cert}
\begin{align}
UB(\mathcal{A}) &= \min_{c \ge 0}\; \sum_{\alpha\in\mathcal{A}} c_\alpha \mu_\alpha
\quad\text{s.t.}\quad \sum_\alpha c_\alpha d^\alpha \ge \prod_i d_i
\;\;\forall d\in\Omega, \label{eq:ub}\\
LB(\mathcal{A}) &= \max_{c \ge 0}\; \sum_\alpha c_\alpha \mu_\alpha
\quad\text{s.t.}\quad \sum_\alpha c_\alpha d^\alpha \le \prod_i d_i
\;\;\forall d\in\Omega. \label{eq:lb}
\end{align}
\end{definition}

\noindent An upper bound is thus \emph{literally} a nonnegativity
certificate for $\sum_\alpha c_\alpha d^\alpha - \prod_i d_i$ on
$\Omega$: classical inequalities --- AM--GM ($abc \le
\frac{a^3+b^3+c^3}{3}$), Young, H\"older --- are the closed-form
members of this cone; the LP selects the best one \emph{adapted to
the observed $\boldsymbol{\mu}$}.

\begin{theorem}[Duality: bounds are extremal alignments]
\label{thm:dual}
The dual of~\eqref{eq:ub} is the \emph{moment LP}
\[
UB(\mathcal{A}) = \max_{\nu \ge 0 \text{ on } \Omega}\; \sum_y \nu(y)\prod_i y_i
\;\;\text{s.t.}\;\; \sum_y \nu(y)\, y^\alpha = \mu_\alpha \;\forall\alpha ,
\]
i.e., the bound equals the largest product moment over all weightings
(alignments) of $\Omega$ consistent with the observed moments; strong
duality holds since the primal is feasible with finite optimum.
Symmetrically, $LB(\mathcal{A})$ is the \emph{best}-case alignment.
\end{theorem}

\noindent This is the semantic reading of a bound: it answers ``how
large can the join be under the worst alignment of degrees consistent
with the data we measured?''  The certificate $c$ is the dual
proof of that answer.

\begin{example}[A three-key star]
\label{ex:three}
Let $Q = R(X,Y) \Join_X S(X,Z)$ on domain $\{x_1,x_2,x_3\}$ with
degree field $d(x_1) = (8,9)$, $d(x_2) = (4,3)$, $d(x_3) = (2,1)$;
hence $|Q| = 72 + 12 + 2 = 86$.  Univariate moments are
$\mu_{(2,0)} = 84$, $\mu_{(0,2)} = 91$, and
$\Omega = \{8,4,2\}\times\{9,3,1\}$.  Minimizing
$84\,c_{(2,0)} + 91\,c_{(0,2)}$ subject to
$c_{(2,0)}a^2 + c_{(0,2)}b^2 \ge ab$ on the nine grid points gives
the Cauchy certificate $(a^2\,t + b^2/t)/2$ at
$t^\star = \sqrt{91/84}$, i.e.\ $UB = \sqrt{84\cdot91} = 87.45$
--- the conjugate H\"older bound.  Now permute the second coordinate
to $(9,3,1) \mapsto (1,3,9)$: \emph{every} univariate moment is
unchanged, so the bound stays $87.45$ while the true join drops to
$8{+}12{+}18 = 38$.  Adding the level-2 moment
$\mu_{(1,1)} = \langle d_1,d_2\rangle$ separates the two instances:
$c_{(1,1)} = 1$ is a certificate on $\Omega$, giving $UB = \mu_{(1,1)}
= |Q|$ exactly in both cases.
\end{example}

\subsection{The H\"older Section}
\label{sec:holder}

Two facts locate the classical bounds inside the framework.

\begin{proposition}[Univariate moments $=$ H\"older, up to the exponent lattice]
\label{prop:holder}
For $\mathcal{A} = \{p\,e_i : p \in P\}$ the certificate LP plays the
role of the conjugate-H\"older bound
$\min_{\sum_i 1/p_i = 1} \prod_i \|\deg_i\|_{p_i}$, with two
differences: (i) the LP's certificates correspond to conjugates
restricted to the lattice $P$ --- a closed-form scan over
\emph{continuous} conjugates can be tighter when the optimum is
non-integral; (ii) dominance is required only on realized degree
values, not all of $\mathbf{R}^m$ --- so the LP is often strictly
tighter than any closed-form point.
\end{proposition}

\begin{proposition}[The entropy LP implies H\"older on pairs]
\label{prop:implied}
In the entropic LP with atom-wise $q$-inequalities~\eqref{eq:qineq},
every connected pair $i,j$ satisfies~\eqref{eq:holderpair}.
\end{proposition}

\begin{proof}
Write $X = A_i\cap A_j$, $U_i = A_i \setminus X$, $U_j = A_j\setminus
X$.  The $q$-inequalities give
$h(X,U_i) \le (1-\frac1p)h(X) + \log\|\deg_i(U_i|X)\|_p$-style bounds
and symmetrically for $j$ with conjugate $q$; elementary
submodularity on the shared separator gives
$h(A_i\cup A_j) \le h(A_i) + h(A_j) - h(X)$.  Summing cancels the
$h(X)$ coefficients since $(1-\tfrac1p)+(1-\tfrac1q)-1 = 0$, leaving
$\log\|\deg_i\|_p + \log\|\deg_j\|_q$.
\end{proof}

\noindent In words: \emph{the entropic LP's implied bound on each pair
is already the H\"older relaxation of the pair join count}; adding
the exact pair moment pins that relaxation to
$\langle \deg_i,\deg_j\rangle$.  The per-pair quantity
\begin{equation}
\mathrm{gap}_{ij} \;=\; \log \frac{\min_{\text{conj.}}\|\deg_i\|_p\|\deg_j\|_q}
{\langle \deg_i, \deg_j \rangle}
\label{eq:gap}
\end{equation}
is therefore a \emph{computable} predictor of what the level-2 moment
buys over level~1 --- the basis of statistic selection under a
budget.

\subsection{The Hierarchy and Exactness}
\label{sec:hier}

For $\mathcal{A}_k = \{\alpha : |\alpha|_1 \le k\}$ the bounds are
monotone, $UB(1) \ge UB(2) \ge \cdots$ --- a Lasserre-style moment
hierarchy~\cite{lasserre2001} over the field.

\begin{proposition}[Top level is exact]
\label{prop:exact}
If $\mathbf{1} = (1,\dots,1) \in \mathcal{A}$ then
$UB(\mathcal{A}) = |Q|$ on star queries, since $\mu_{\mathbf{1}}$
\emph{is} the product moment.  Hence $m$-atom stars need the $m$-th
moment; on STATS, all $m{=}3$ stars are exact at level~3.
\end{proposition}

Support information in $\Omega$ tightens this further:

\begin{lemma}[Key collapse]
\label{lem:collapse}
If $\mathrm{supp}(\deg_i) \subseteq \{0,1\}$ (a uniqueness constraint
on the separator), then $\alpha = \mathbf{1} - e_i$ is a feasible
monomial certificate and
$UB \le \mu_{\mathbf{1}-e_i} = |Q| +
\sum_{x\in(\cap_{j\ne i}\mathrm{supp}_j)\setminus \mathrm{supp}_i}
\prod_{j\ne i} d_j(x)$.  Exactness at level $m{-}1$ therefore needs
both $d_i \equiv 1$ \emph{and} $\mathrm{supp}_i \supseteq
\cap_{j\ne i}\mathrm{supp}_j$.  On STATS $m{=}4$ stars, all six
exactly-estimated queries carry a $\max\deg{=}1$ atom; on the single
query with a collapsed-but-not-covering atom the bound overshoots by
exactly the predicted uncovered mass.
\end{lemma}

\subsection{The Lower-Bound Wing}
\label{sec:lbwing}

The framework is symmetric: $LB(\mathcal{A})$ of~\eqref{eq:lb} is the
best-case alignment.  Here the polynomial cone honestly fails: on
stars with $m\ge 3$ the best monomial certificate under $\prod d_i$ is
$0$, because the inclusion--exclusion term
$\max(0,\sum_i d_i - (m{-}1))$ is piecewise-linear, not polynomial.
The matching basis is the family of \emph{hinge functions} of the
form $d \mapsto \max(0,\langle w,d\rangle - t)$, whose field
expectations are the
inclusion--exclusion moments; bounds remain one object --- extremal
moments over alignments --- with certificates as conic combinations
of a chosen basis.

\subsection{Scalability: Small LPs and Separation}
\label{sec:sep}

The certificate LP has $|\mathcal{A}|$ variables and $|\Omega|$ rows;
$\Omega$ is exponential only in $m$ (atoms on the separator), not in
the $n$ attributes of the query --- the $2^n$ blow-up of the entropic
LP disappears.  When $|\Omega| = \prod_i |\mathrm{supp}_i|$ is large
we solve~\eqref{eq:ub} by \emph{separation} (cutting planes): solve on
a seed domain, query an oracle for the most-violated grid point, add
it, repeat.  Convergence returns a certificate dominating the
\emph{entire} product --- soundness is proved, not sampled.
